\section{The closure theorem}

\begin{theorem}[Encounter-orientation receipt tower]
Let \(O\cong C_4\) be a visible orientation cycle with a deterministic
binary two-sheet lift of nontrivial circuit holonomy. Let \(H\) be its
lifted successor. If \(H=AB\) for two involutions \(A\) and \(B\),
then, up to fiber-preserving gauge and labeled group isomorphism:
\begin{enumerate}
\item the lifted orientation cycle is \(\langle H\rangle\cong C_8\);
\item the generated completion is \(\langle A,B\rangle\cong D_{16}\);
\item the induced encounter quotient is \(E\cong C_2\times C_2\);
\item the extension receipt is \(c=[A,B]=H^2\);
\item the binary return receipt is \(\tau=c^2=H^4\).
\end{enumerate}
The encounter quotient and orientation quotient are nonisomorphic
four-state groups.
\end{theorem}

\begin{proof}
Nontrivial binary holonomy forces the connected eight-cycle by the
binary lift classification, so \(H\) has order eight. Two involutions
whose product has order eight generate the dihedral group of order
sixteen. The commutator calculation gives \([A,B]=H^2\), and squaring
gives \([A,B]^2=H^4\), the unique nonidentity element in the deck
kernel of the \(C_8\to C_4\) projection. Mapping \(A\) and \(B\) to
the independent involutions of \(V_4\) gives the encounter quotient.
The distinct element-order profiles prove nonidentification.
\end{proof}

The theorem supplies the exact chain
\[
\text{ordered encounter}
\longrightarrow [A,B]
\longrightarrow [A,B]^2
\longrightarrow \text{binary holonomy}
\longrightarrow C_8
\longrightarrow D_{16}.
\]
